\documentclass[11pt]{article}

\usepackage[letterpaper,margin=1in]{geometry}
\usepackage{amsmath,amssymb,amsthm}
\usepackage[colorlinks=true,allcolors=blue]{hyperref}

\newtheorem{prop}{Proposition}
\newtheorem{lem}{Lemma}
\newtheorem{defi}{Definition}
\newtheorem{rmk}{Remark}

\newcommand{\cA}{\mathcal A}
\newcommand{\cD}{\mathcal D}
\newcommand{\cI}{\mathcal I}
\newcommand{\cT}{\mathcal T}
\newcommand{\cU}{\mathcal U}
\newcommand{\E}{\mathbb E}
\newcommand{\Pbb}{\mathbb P}
\newcommand{\bp}{\mathbf p}
\newcommand{\bn}{\mathbf n}
\newcommand{\ileft}{\mathrm I_{\rm left}}
\newcommand{\iright}{\mathrm I_{\rm right}}
\newcommand{\stable}{\mathrm{stable}}

\newcommand{\mscomment}[1]{\textup{\textsf{\color{red}[#1]}}}
\newcommand{\jwcomment}[1]{\textup{\textsf{\color{purple}[#1]}}}
\newcommand{\jscomment}[1]{\textup{\textsf{\color{blue}[#1]}}}

\title{Time-Inhomogeneous Riffle Shuffles: First-Moment Section}
\author{}
\date{}

\begin{document}
\maketitle

\jscomment{I think the following are a general problems. 1. We should ask ai to add more "recall the definition of ..." part. 2.AI prefers to use vague phrase instead of referring to Lemma x.x or (x.x). 3. AI likes creating some new phrase without explaining its meaning rigorously. We should ask AI to avoid using any new concept without explaining its meaning rigorously in math sense before.}

\section{Time-Inhomogeneous First-Moment Estimate}

This section isolates the first-moment estimate needed in the
time-inhomogeneous upper bound. The proof follows the homogeneous
first-moment argument, with the single measure \(\mu\) replaced by a
time-profile \((\mu_t)\) and with all numerical hypotheses written as
space-time integral inequalities.

Fix \(k\ge2\), write \([k]_0=\{0,\ldots,k-1\}\), and set
\[
    \cD_k=\left\{\bp=(p_0,\ldots,p_{k-1})\in[0,1]^k:\ \sum_{\ell=0}^{k-1}p_\ell=1\right\}.
\]
Fix \(T<\infty\), let \(N\) be a positive integer, set
\(K=\lfloor T\log N\rfloor\), and fix deterministic pile sizes, indexed
by \(1\le j\le K\),
\[
    \bn^{(j)}=(n_0^{(j)},\ldots,n_{k-1}^{(j)}),\qquad
    n_\ell^{(j)}\in\mathbb Z_{\ge0},\qquad
    \sum_{\ell=0}^{k-1}n_\ell^{(j)}=N .
\]
Write \(\bp^{(j)}=\bn^{(j)}/N\in\cD_k\). Let \((\mu_t)_{0\le t\le T}\) be a Borel probability kernel on \(\cD_k\). For \(\bp\in\cD_k\), define
\[
    p_{\max}(\bp)=\max_\ell p_\ell,
    \qquad
    \psi_{\bp}(\theta)=-\log\sum_{\ell=0}^{k-1}p_\ell^\theta .
\]
Set
\[
    a(t)=\E_{\mu_t}\log\frac1{p_{\max}(\bp)},\qquad
    \psi_t(\theta)=\E_{\mu_t}\psi_{\bp}(\theta),
\]
and
\[
    A(r,s)=\int_r^s a(t)\,dt,\qquad
    B_\theta(r,s)=\int_r^s\psi_t(\theta)\,dt .
\]

A word \(x=(x[1],\ldots,x[M])\in[k]_0^M\), \(0\le M\le K\),
is legal if \(n_{x[j]}^{(j)}>0\) for every \(1\le j\le M\). Let
\(\cT\) be the tree of legal words of depth at most \(K\), ordered by
prefix. The random shuffle matrix is sampled column by column: in column
\(j\), the \(N\) entries are a uniformly random permutation of the
multiset containing \(n_\ell^{(j)}\) copies of \(\ell\). Sorting the
resulting rows lexicographically gives intervals \(\cI(B)\) for prefix
cylinders \(B\). If \(x\in\cT\), then \(M_x=|x|\), \(B_x\) denotes the
set of \jscomment{"full legal strings"-->legal strings with length $K$?} extending \(x\), and \(\cI(B_x)\) is its random
lexicographic interval. If \(\cU\) is a prefix \jscomment{->antichain? Is it necessary to introduce this concept here? If so, add a footnote explain this.}, set
\(B_\cU=\bigcup_{x\in\cU}B_x\) and
\(\cI(\cU)=\cI(B_\cU)\). For a non-root word \(x\), \(\mathcal F x\)
denotes its parent.

The shuffle graph \(G\) has vertex set \([N]\) after the lexicographic
sort and edge set
\[
    E(G)=\{(m,m+1):\text{ the sorted rows at positions }m,m+1
    \text{ are identical}\}.
\]
For a prefix \(x\), \(G_{B_x}\) is the subgraph containing those edges of
\(G\) whose two endpoint rows lie in \(B_x\). For two independent shuffle
graphs \(G,G'\), write
\[
    E(G,G')=E(G)\cap E(G')
\]
for their common adjacent edges.

\begin{defi}[Almost time-profile-like cuts]
Given a finite measurable partition \(\cD_k=D_0\cup\cdots\cup D_m\), the deterministic cuts are \((\rho,\varphi)\)-almost-\(\mu_\cdot\)-like on \([0,T]\) if, for every grid interval
\[
    I=[a,\min(a+\rho,T)]\subset[0,T],
    \qquad a\in\rho\mathbb Z,
\]
with positive length, and every \(D_i\),
\[
    \left|
    \frac1{\log N}\#\{1\le j\le K:\ j/\log N\in I,\ \bp^{(j)}\in D_i\}
    -\int_I\mu_t(D_i)\,dt
    \right|\le \varphi |I|+\frac{1}{\log N}.
\]
\end{defi}

We say the cuts satisfy the local \((\chi,\gamma,\rho)\)-good-shuffle
condition on an interval \([r,T]\), with fixed \(\gamma>0\), if every
macroscopic interval \(I\subset[r,T]\) of length \(\rho\) contains at
least \(\gamma\rho\log N\) indices \(1\le j\le K\) with
\(j/\log N\in I\) for which at least two piles have size at least
\(\chi N\).

The parameters below are chosen in the order
\[
    k,T,\eta,\gamma\ \text{fixed},\qquad
    \chi\downarrow0,\quad
    \alpha_\chi=(\log(1/\chi))^{-1/2},\quad
    \Delta,\rho,\varphi,\delta\downarrow0,\qquad
    N\to\infty .
\]
The choices are nested in the displayed order; in particular all constants
may depend on the parameters to their left. All \(o(1)\) terms are with
this final limit, uniformly over the finitely many retained profile
classes after the truncations below. Statements of the form
\(N^{o_\chi(1)}\) mean that for every \(\varepsilon>0\), after choosing
\(\chi\) sufficiently small and then \(N\) sufficiently large, the term
is at most \(N^\varepsilon\).

\begin{defi}[\(\chi\)-regular partitions]
A finite Borel partition \(\cD_k=D_0\cup\cdots\cup D_m\), with
representatives \(\bp^{(i)}\in D_i\), is called \(\chi\)-regular if:
\[
    \sup_{\bp,\bp'\in D_i}\sup_{\theta\in[2,4]}
    |\psi_{\bp}(\theta)-\psi_{\bp'}(\theta)|\le C\chi,
    \qquad
    \sup_{\bp,\bp'\in D_i}
    \left|\log\frac1{p_{\max}(\bp)}-\log\frac1{p_{\max}(\bp')}\right|
    \le C\chi,
\]
and, for every coordinate \(\ell\),
\[
    p_\ell^{(i)}>2\chi,\ \bp\in D_i
    \quad\Longrightarrow\quad
    p_\ell>\chi,\qquad
    \left|\log\frac1{p_\ell}-\log\frac1{p_\ell^{(i)}}\right|\le C\chi,
\]
while
\[
    p_\ell^{(i)}\le2\chi,\ \bp\in D_i
    \quad\Longrightarrow\quad p_\ell\le3\chi .
\]
For each fixed \(\chi\), such a partition is obtained by subdividing the
simplex into finitely many Borel coordinate cells of diameter
\(O(\chi^2)\), with additional threshold cuts at the coordinates
\(\chi,2\chi,3\chi\). In what follows, this regular partition is chosen
first, and the almost-\(\mu_\cdot\)-like hypothesis is imposed for this
chosen partition.
\end{defi}

\jscomment{I am not sure. Would it be better to first give the construction of $D_i$, and then claim that (and even prove a lemma maybe) $D_i$ has the properties listed in the definition above?}

\begin{defi}[Ambient block and size exponent]
An ambient block is a fixed legal prefix block \(B_w\) with
\(M_w=|w|\), \(M_w/\log N=r+o(1)\), and expected size
\[
    \E|\cI(B_w)|=N^{b+o(1)}.
\]
\jscomment{1. Is the ambient block should be defined for a pair $(r,b)$? 2. Should we take "ambient block" as a property of $B_w$? If so, the first sentence of the definition should be rewritten. 3. The notation $o(1)$ should be related to some limiting process, but here in the definition only one prefix $w$ is considered.}
The whole deck corresponds to \(r=0\), \(b=1\), and \(w=\varnothing\). For any legal extension prefix \(x\supseteq w\), with absolute length \(M_x\), define
\[
    t_x=\frac{M_x}{\log N},
\]
and
\[
    c_L^{r,b}(x)
    =
    b-\frac1{\log N}\sum_{M_w<j\le M_x}
    \log\frac{N}{n_{x[j]}^{(j)}} .
\]
Thus \(c_L^{r,b}(x)=\log_N\E|\cI(B_x)|+o(1)\) inside the ambient block.
When an argument is carried out inside \(B_w\), interval endpoints and
edge locations are recentered at the left endpoint of \(\cI(B_w)\);
translating the whole ambient block is irrelevant for internal edge
counts. \jscomment{It seems that the sentence before should not appear in a definition...}
Also define the tail edge exponent, for every legal prefix \(x\supseteq w\), by
\[
    c_E(x)=-\frac1{\log N}\sum_{M_x<j\le K}\psi_{\bp^{(j)}}(2).
\]
All quantities carrying only the superscript \((r,b)\) in this section are
understood to be relative to this fixed ambient prefix \(w\).
\jscomment{Is $(r,b)$ completely determined by $w$? If not, would it be better to also write the dependence on $w$ explicitly?}
\end{defi}

\begin{defi}[Stable blocks]
\jscomment{This definition is about 2-pages long, and I think we should decompose it into several definitions...}
Fix \(0<\chi<1/k\) and a large constant \(\mathsf C\). For each shuffle time \(j\), the sets below are nonempty, and we set
\[
    \ileft^{(j)}=\inf\{\ell:\ n_\ell^{(j)}/N\ge\chi\},
    \qquad
    \iright^{(j)}=\sup\{\ell:\ n_\ell^{(j)}/N\ge\chi\}.
\]
The left boundary-prefix family inside \(B_w\) is
\[
    \begin{aligned}
    \cA_{\rm left}^{r,b}
    &=
    \left[\bigcup_{M_w<M\le K}
    \left\{x\in\cT:
    \begin{array}{l}
    x\supseteq w,\ |x|=M,\ x[M]>\ileft^{(M)},\\
    x[j]\le\ileft^{(j)}\text{ for }M_w<j<M
    \end{array}\right\}\right]
    \\
    &\quad\cup
    \left\{x\in\cT:
    \begin{array}{l}
    x\supseteq w,\ |x|=K,\\
    x[j]\le\ileft^{(j)}\text{ for all }M_w<j\le K
    \end{array}\right\}.
    \end{aligned}
 \]
Similarly, the right boundary-prefix family is
\[
    \begin{aligned}
    \cA_{\rm right}^{r,b}
    &=
    \left[\bigcup_{M_w<M\le K}
    \left\{x\in\cT:
    \begin{array}{l}
    x\supseteq w,\ |x|=M,\ x[M]<\iright^{(M)},\\
    x[j]\ge\iright^{(j)}\text{ for }M_w<j<M
    \end{array}\right\}\right]
    \\
    &\quad\cup
    \left\{x\in\cT:
    \begin{array}{l}
    x\supseteq w,\ |x|=K,\\
    x[j]\ge\iright^{(j)}\text{ for all }M_w<j\le K
    \end{array}\right\}.
    \end{aligned}
 \]

For every full legal string \(\omega\in B_w\), let \(L(\omega)\) be the
shortest prefix of \(\omega\) belonging to
\(\cA_{\rm left}^{r,b}\), and let \(R(\omega)\) be the shortest prefix
belonging to \(\cA_{\rm right}^{r,b}\). These prefixes are \jscomment{-->unique? If shortest, it is of course unique as a prefix.} because
the definitions stop at \jscomment{-->the first genuine crossing?}of the left or right
large-pile boundary, with the terminal word at time \(K\) used when no
crossing occurs.\jscomment{Another explanatory sentence that should not appear in a definition...} Thus every full string in \(B_w\) is assigned a unique
boundary pair \((L(\omega),R(\omega))\).

The admissible boundary families are the truncations
\[
    \widetilde\cA_{\rm left}^{r,b}
    =
    \left\{y\in\cA_{\rm left}^{r,b}:\
    \#\{M_w<j\le M_y:\ y[j]<\ileft^{(j)}\}
    \le \frac{\mathsf C}{\log(1/\chi)}\log N\right\},
\]
\[
    \widetilde\cA_{\rm right}^{r,b}
    =
    \left\{z\in\cA_{\rm right}^{r,b}:\
    \#\{M_w<j\le M_z:\ z[j]>\iright^{(j)}\}
    \le \frac{\mathsf C}{\log(1/\chi)}\log N\right\}.
\]

For \(y\in\widetilde\cA_{\rm left}^{r,b}\) and \(z\in\widetilde\cA_{\rm right}^{r,b}\), define
\[
    A_{y,z}^{r,b}=\{x\in\cT:\ x\supseteq w,\ x\text{ has both }y\text{ and }z\text{ as prefixes}\}.
\]
When this is nonempty, one of \(y,z\) extends the other. Write \(y\vee z\) for the longer of \(y,z\), and set
\[
    s=s_{y,z}=t_{y\vee z}.
\]
For \(x\in A_{y,z}^{r,b}\), let \(M_x,M_y,M_z\) be the absolute lengths
of \(x,y,z\), and put \(M_\vee=|y\vee z|\). Define the left
lexicographic pieces inside the ambient block by
\[
    P_{x,i}^{r,b}
    =
    \{\sigma\in\cT:\ |\sigma|=M_x,\ \sigma \text{ extends }w,\ 
    \sigma[j]=x[j]\text{ for }M_w<j<i,\ \sigma[i]<x[i]\},
\]
for \(M_w<i\le M_x\). Similarly define the right pieces
\[
    Q_{x,i}^{r,b}
    =
    \{\sigma\in\cT:\ |\sigma|=M_x,\ \sigma \text{ extends }w,\ 
    \sigma[j]=x[j]\text{ for }M_w<j<i,\ \sigma[i]>x[i]\}.
\]
All expectations below are for the random strings generated inside the ambient block \(B_w\), and for a family of prefixes we use the interval convention above. Set
\[
    \mathrm{Fluc}_{\rm left}^{r,b}(x;y)
    =
    \left(\max_{M_w<i\le M_\vee}\E |\cI(P_{x,i}^{r,b})|\right)^{1/2},
\]
\[
    \mathrm{Fluc}_{\rm right}^{r,b}(x;z)
    =
    \left(\max_{M_w<i\le M_\vee}\E |\cI(Q_{x,i}^{r,b})|\right)^{1/2},
\]
and
\[
    c_F^{r,b}(x;y,z)
    =
    \min\left(
    \log_N \mathrm{Fluc}_{\rm left}^{r,b}(x;y),
    \log_N \mathrm{Fluc}_{\rm right}^{r,b}(x;z)
    \right).
\]
\jscomment{I am just curious: In our time-homogeneous paper, the maximum in the definition of $\mathrm{Fluc}$ is taken up to $M_y$ or $M_z$, but here is taken up to $M_\vee=|y\vee z|$. Are they equivalent? If not, why this change is made?}
The maximum over an empty index set is \(0\), and \(\log_N0=-\infty\).
Because \(x\) has both \(y\) and \(z\) as prefixes, the quantities above
depend only on the boundary pair \((y,z)\), not on the continuation beyond
\(y\vee z\). We use the representative value
\(c_F^{r,b}(y\vee z;y,z)\) below.\jscomment{"representative" in the sense that \(c_F^{r,b}(y\vee z;y,z)=c_F^{r,b}(x;y,z)\)?}

Call a pair \((y,z)\) terminal if \(M_y=K\) or \(M_z=K\). A pair is
retained if \(A_{y,z}^{r,b}\ne\varnothing\) and either it is terminal or
\(c_F^{r,b}(y\vee z;y,z)>-\infty\). A nonterminal retained pair is called
fluctuation-regular when
\[
    c_F^{r,b}(\xi;y,z)\ge \frac12c_L^{r,b}(\xi)-o(1)
\]
holds uniformly for all continuations \(\xi\in A_{y,z}^{r,b}\). This is
proved in the profile-estimates lemma below from the genuine left and
right boundary crossings:
at the crossing time for \(y\) there is at least \(\chi N\) mass on the
left side, and at the crossing time for \(z\) there is at least
\(\chi N\) mass on the right side. Terminal retained pairs are not
discarded; their stable family is the singleton \(\{y\vee z\}\) and they
are handled by the terminal/tail estimate. \jscomment{Another explanatory sentences below the displayed equation that should not appear in a definition...}

For fixed \(\delta>0\), define the first-crossing and terminal-fallback sets
\[
    \Gamma_{y,z}^{r,b}
    =
    \{x\in A_{y,z}^{r,b}:\
    c_L^{r,b}(x)-c_F^{r,b}(y\vee z;y,z)<\delta,\
    c_L^{r,b}(\mathcal F x)-c_F^{r,b}(y\vee z;y,z)\ge\delta\},
\]
\[
    \Lambda_{y,z}^{r,b}
    =
    \{x\in A_{y,z}^{r,b}:\ |x|=K,\
    c_L^{r,b}(x)-c_F^{r,b}(y\vee z;y,z)\ge\delta\}.
\]
Here \(\mathcal F x\) is the parent of \(x\), used only in \(\Gamma_{y,z}^{r,b}\). Now set
\[
    A_{y,z}^{\stable,r,b}
    =
    \begin{cases}
    \{y\vee z\},
    & \text{\((y,z)\) terminal},\\[4pt]
    \{y\vee z\},
    & c_L^{r,b}(y\vee z)-c_F^{r,b}(y\vee z;y,z)<\delta,\\[4pt]
    \Gamma_{y,z}^{r,b}\cup\Lambda_{y,z}^{r,b},
    & c_L^{r,b}(y\vee z)-c_F^{r,b}(y\vee z;y,z)\ge\delta,
    \end{cases}
\]
%\jscomment{Could you explain some difference between this definition and the definiton in the time-homogeneous case?1. There is no the $\Lambda$ set in the definition therein. 2. There is not the case that \((y,z)\) terminal.}
For stable \(x\), set \(u=t_x\). Thus every stable block has times
\[
    r\le s\le u\le T.
\]
The set \(\Lambda_{y,z}^{r,b}\) is included only to make the stopping
rule total before applying the tail-size margin; in the first-moment
application, the tail-size part of \(\mathsf{FM}(r,b,T,\eta)\) excludes
nonterminal fallback blocks, while terminal singleton pairs are estimated
directly by the same tail-size margin. \jscomment{Maybe it is better to make these explanations as a footnote.}
\end{defi}

\begin{defi}[Space-time digit profile]
Fix the \(\chi\)-regular partition
\(\cD_k=D_0\cup\cdots\cup D_m\) from the almost-like condition, with
representatives \(\bp^{(i)}\in D_i\), and choose a time mesh
\[
    0=t_0<t_1<\cdots<t_R=T,\qquad J_q=[t_q,t_{q+1})\ (q<R-1),\quad J_{R-1}=[t_{R-1},T]
\]
with mesh size at most \(\Delta\). Thus \(q\in\{0,\ldots,R-1\}\) is only the time-cell index. Let \((y,z)\) be a retained boundary pair, let \(x\in A_{y,z}^{\stable,r,b}\), and put \(s=s_{y,z}\), \(u=t_x\).

For a shuffle time \(j\) with \(s\log N<j\le u\log N\), let \(q(j)\) and \(i(j)\) be the unique indices such that
\[
    j/\log N\in J_{q(j)},\qquad \bp^{(j)}\in D_{i(j)}.
\]
Call the pair \((j,x[j])\) good if \(p_{x[j]}^{(i(j))}>2\chi\), and exceptional otherwise.

Define
\[
    c_\ell^{(q,i)}(x)
    =
    \frac1{\log N}
    \#\{j:\ s\log N<j\le u\log N,\ q(j)=q,\ i(j)=i,\ x[j]=\ell,\ (j,x[j])\text{ is good}\}.
\]
The full profile consists of the counts \(c_\ell^{(q,i)}(x)\), together with the exceptional set
\[
    \mathsf Q_x=\{j:\ s\log N<j\le u\log N,\ (j,x[j])\text{ is exceptional}\}.
\]
Define
\[
    S_x=\sum_{q=0}^{R-1}\sum_{i=0}^m\sum_{\ell=0}^{k-1}
    c_\ell^{(q,i)}(x)\log\frac1{p_\ell^{(i)}},
\]
\[
    Q_x=\frac1{\log N}\sum_{j\in\mathsf Q_x}
    \log\frac{N}{n_{x[j]}^{(j)}}\ge0,
    \qquad
    R_x=\frac{|\mathsf Q_x|}{\log N}.
\]
For \(c_{\rm tot}^{(q,i)}=\sum_\ell c_\ell^{(q,i)}\), define
\[
    H_x
    =
    \sum_{q,i:\ c_{\rm tot}^{(q,i)}>0}
    c_{\rm tot}^{(q,i)}
    \sum_{\ell:\ c_\ell^{(q,i)}>0}
    \frac{c_\ell^{(q,i)}}{c_{\rm tot}^{(q,i)}}
    \log\frac{c_{\rm tot}^{(q,i)}}{c_\ell^{(q,i)}} .
\]
\end{defi}

\begin{defi}[Numerical first-moment margin]
For an ambient block \((r,b)\), set
\[
    \beta_{r,b}(s)=b-A(r,s).
\]
For \(r\le s\le u\le T\), define
\[
    \mathcal N_1(r,s,u;b)
    =
    \inf_{\theta\in[3,4]}
    \left\{\frac{\theta+3}{2}\beta_{r,b}(s)-B_\theta(s,u)-2B_2(u,T)\right\},
\]
\[
    \mathcal N_2(r,s,u;b)
    =
    \inf_{\theta\in[3,4]}
    \left\{\frac{\theta-1}{2}\beta_{r,b}(s)-B_\theta(s,u)-2B_2(u,T)\right\},
\]
and
\[
    \mathcal N_3(r,s;b)=\beta_{r,b}(s)-B_2(s,T).
\]
We say \(\mathsf{FM}(r,b,T,\eta)\) holds if
\[
    \beta_{r,b}(T)\le -\eta,
\]
and
\[
    \mathcal N_1(r,s,u;b)\le-\eta,\qquad
    \mathcal N_2(r,s,u;b)\le-\eta
\]
for all \(r\le s\le u\le T\), and
\[
    \mathcal N_3(r,s;b)\le-\eta
\]
for all \(r\le s\le T\). Write \(\mathsf{FM}(T,\eta)=\mathsf{FM}(0,1,T,\eta)\).
\end{defi}

\jscomment{Maybe we can add a remark here about which terms are the leading term, i.e.\ giving the constant in the cutoff time.}

\begin{lem}[Stable partition]
For every retained boundary pair \((y,z)\), the stable blocks \(B_x\),
\(x\in A_{y,z}^{\stable,r,b}\), partition the strings in
\(B_{y\vee z}\). The intervals \(\cI(B_x)\) partition
\(\cI(B_{y\vee z})\).
\end{lem}

\begin{proof}
Fix a retained pair \((y,z)\), and write \(v=y\vee z\). If
\((y,z)\) is terminal, then by definition
\(A_{y,z}^{\stable,r,b}=\{v\}\), so the assertion is immediate. If
\(c_L^{r,b}(v)-c_F^{r,b}(v;y,z)<\delta\), then by definition
\(A_{y,z}^{\stable,r,b}=\{v\}\), so the assertion is immediate.

Otherwise consider a full legal string \(\omega\in B_v\). Along the
unique path from \(v\) to \(\omega\), the quantity \(c_L^{r,b}(\cdot)\)
is nonincreasing, while \(c_F^{r,b}(\cdot;y,z)\) is fixed because the
fluctuation pieces only use coordinates up to the boundary lengths
\(M_y\) and \(M_z\). If a first descendant \(x\) of length at most \(K\)
satisfies
\[
    c_L^{r,b}(x)-c_F^{r,b}(v;y,z)<\delta,
\]
then \(x\in\Gamma_{y,z}^{r,b}\); this includes the possibility
\(|x|=K\). If no such first crossing occurs, the terminal word of
\(\omega\) lies in \(\Lambda_{y,z}^{r,b}\). Thus every full string in
\(B_v\) has a unique candidate stable prefix.

The family \(\Gamma_{y,z}^{r,b}\cup\Lambda_{y,z}^{r,b}\) is an
antichain. Indeed, if \(x\in\Gamma_{y,z}^{r,b}\), then every strict
descendant already lies after the first crossing and cannot satisfy the
parent condition in \(\Gamma\), while it also cannot lie in
\(\Lambda\), whose elements have no crossing before the terminal time.
Distinct \(\Lambda\)-elements are leaves. Therefore the stable prefix
cylinders are disjoint. Since prefix cylinders are contiguous in
lexicographic order, their images are disjoint contiguous subintervals
whose union is \(\cI(B_v)\).
\end{proof}

\begin{lem}[Global boundary decomposition]
Before the admissibility truncation, the maps \(L(\omega)\) and
\(R(\omega)\) assign every full legal string \(\omega\in B_w\) to a
unique boundary pair. Consequently
\[
    B_w=
    \bigsqcup_{(y,z)}
    B_{\{\,\omega:\ L(\omega)=y,\ R(\omega)=z\,\}},
\]
where the disjoint union is over boundary pairs from
\(\cA_{\rm left}^{r,b}\times\cA_{\rm right}^{r,b}\). For every admissible
retained pair \((y,z)\), the stable cylinders
\(\{B_x:x\in A_{y,z}^{\stable,r,b}\}\) partition the corresponding
subfamily \(B_{y\vee z}\). The only strings not represented by retained
admissible stable cylinders are those for which \(L(\omega)\) or
\(R(\omega)\) lies outside the admissible truncation.
\end{lem}

\jscomment{Could you also explain or recall what do you mean by "admissible truncation"?}

\begin{proof}
Uniqueness of \(L(\omega)\) and \(R(\omega)\) follows from first-crossing
minimality, with the terminal prefix used when no crossing occurs. This
gives the first disjoint decomposition. If a boundary pair is admissible
and retained, the stable partition lemma partitions \(B_{y\vee z}\) into its stable
cylinders. If a string is not covered after the admissibility truncation,
its unique boundary pair has at least one component removed by the
truncation. Empty pairs contribute no strings.
\end{proof}

\begin{lem}[Regular partition and interval blocking]
Assume the cuts are \((\rho,\varphi)\)-almost-\(\mu_\cdot\)-like for the
chosen \(\chi\)-regular partition. Then the representative replacements
used in the profile estimates have error \(O(\chi)\), good digits have
actual pile mass at least \(\chi\), and exceptional digits have actual
pile mass at most \(3\chi+o(1)\). Moreover, the \(\rho\)-grid
almost-like condition implies the same cell-count estimates on arbitrary
subintervals of \([0,T]\), with error \(O(\rho+\varphi)+o(1)\).
\end{lem}

\jscomment{
1. "The representative replacements used in the profile estimates", "\(\rho\)-grid
almost-like condition" and "same cell-count estimates on arbitrary
subintervals of \([0,T]\)" are not a rigorous math notation/expression should appear in a statement of lemma.
Try to add some reference.
2. Recall the definition of "good" and "exceptional" from Definition 5.
}

\begin{proof}
The representative replacement statements are precisely the defining
properties of a \(\chi\)-regular partition. If
\(\bp^{(j)}\in D_i\) and \(p_\ell^{(i)}>2\chi\), then
\(n_\ell^{(j)}/N=p_\ell^{(j)}>\chi\) and
\[
    \left|
    \log\frac{N}{n_\ell^{(j)}}-\log\frac1{p_\ell^{(i)}}
    \right|\le C\chi .
\]
If \(p_\ell^{(i)}\le2\chi\), then
\(n_\ell^{(j)}/N\le3\chi\). Thus exceptional digits carry an actual
small-pile factor \(O(\chi)\), while good digits are in the controlled
logarithmic region.

For interval blocking, cover any interval \([a,b]\subset[0,T]\) by the
maximal union of half-open full \(\rho\)-grid intervals it contains,
using the closed final interval only at \(T\). The omitted left and right
boundary pieces have total length at most \(2\rho\). The floor errors in
converting times to indices are \(O(1/\log N)\) per grid block; since
\(\rho\) is fixed before \(N\to\infty\), their total contribution is
o(1). Applying the almost-like condition on the full grid intervals and
summing over them gives, uniformly in \(i\),
\[
    \frac1{\log N}\#\{j:\ j/\log N\in[a,b],\ \bp^{(j)}\in D_i\}
    =
    \int_a^b\mu_t(D_i)\,dt+O(\rho+\varphi)+o(1).
\]
\end{proof}

\begin{lem}[Profile estimates]
Uniformly over retained boundary pairs and stable
\(x\in A_{y,z}^{\stable,r,b}\),
\[
    c_L^{r,b}(x)=c_L^{r,b}(y\vee z)-S_x-Q_x
    +O(\Delta+\chi+\rho+\varphi)+o(1),
\]
\[
    c_E(x)=-B_2(u,T)+O(\Delta+\chi+\rho+\varphi)+o(1),
\]
\[
    c_L^{r,b}(y\vee z)\le
    \beta_{r,b}(s)+O(\Delta+\chi+\rho+\varphi)+o(1).
\]
With \(\alpha_\chi=(\log(1/\chi))^{-1/2}\), every continuation
\(\xi\in A_{y,z}^{r,b}\) of a nonterminal fluctuation-regular retained
pair satisfies
\[
    c_F^{r,b}(\xi;y,z)\ge \frac12 c_L^{r,b}(\xi)-o(1).
\]
\jscomment{What is the difference between this statement and the definition of "nonterminal fluctuation-regular retained
pair" except the "continuation"?}

If \(x\in A_{y,z}^{\stable,r,b}\) is in the non-singleton first-crossing
case, then
\[
    S_x\le c_L^{r,b}(y\vee z)-c_F^{r,b}(y\vee z;y,z)
    +O(\Delta+\chi+\rho+\varphi+\alpha_\chi)+o(1).
\]
\jscomment{"After discarding profiles with \(R_x>\alpha_\chi\)" Do you mean by this the discarding in the definition of $H_x$? If so, it seems unnecessary to mention this again here.}, for every
\(\theta\in[3,4]\),
\[
    H_x\le \theta S_x-B_\theta(s,u)
    +O(\Delta+\chi+\rho+\varphi+\alpha_\chi)+o(1).
\]
\end{lem}

\jscomment{Write or recall that $u=t_x$ in this lemma.}

\begin{proof}
Write \(v=y\vee z\). The identity for \(c_L^{r,b}(x)\) is the definition
of \(c_L^{r,b}\) split at time \(s=t_v\).
On good positions, the regular
partition lemma replaces the actual logarithm by the representative
logarithm with total error \(O(\chi)\); exceptional positions are kept
exactly in \(Q_x\). Hence
\[
    c_L^{r,b}(x)=c_L^{r,b}(v)-S_x-Q_x+O(\chi)+o(1),
\]
and we keep the larger displayed error to match the tail and entropy
bookkeeping.

\jscomment{I think we should ask gpt to use exact reference like Lemma... or equation (...) instead of expression like "regular partition and interval-blocking lemma".}

For \(c_E(x)\), group the tail times \(u\log N<j\le K\) by time cells
and simplex cells. The regular partition and interval-blocking lemma
give
\[
    -\frac1{\log N}\sum_{M_x<j\le K}\psi_{\bp^{(j)}}(2)
    =
    -\int_u^T\psi_t(2)\,dt
    +O(\Delta+\chi+\rho+\varphi)+o(1).
\]
The same argument with \(\log(1/p_{\max})\) gives the maximal-pile bound
for \(v\):
\[
    \frac1{\log N}\sum_{M_w<j\le M_v}
    \log\frac1{p_{\max}(\bp^{(j)})}
    =
    A(r,s)+O(\Delta+\chi+\rho+\varphi)+o(1).
\]
Since \(n_{v[j]}^{(j)}/N\le p_{\max}(\bp^{(j)})\), this yields the
displayed upper bound on \(c_L^{r,b}(v)\).

\jscomment{I believe the two displayed equation is correct, but shall we provide a more detailed proof for them?}

We next record the fluctuation comparison. Suppose the left boundary
prefix \(y\) has a genuine crossing before time \(K\), and write
\(i=M_y\). Then \(y[i]>\ileft^{(i)}\), so
\[
    \sum_{\ell<y[i]}n_\ell^{(i)}\ge n_{\ileft^{(i)}}^{(i)}\ge\chi N .
\]
For every continuation \(\xi\in A_{y,z}^{r,b}\), the side piece
\(P_{\xi,i}^{r,b}\) contains the subfamily obtained from \(B_\xi\) by
replacing the \(i\)-th digit \(y[i]\) by a digit \(<y[i]\) and keeping all
later digits fixed. 
Hence, inside the ambient block \jscomment{which ambient block},
\[
    \E|\cI(P_{\xi,i}^{r,b})|
    \ge
    \E|\cI(B_\xi)|
    \frac{\sum_{\ell<y[i]}n_\ell^{(i)}}{n_{y[i]}^{(i)}}
    \ge
    \chi\,\E|\cI(B_\xi)|.
\]
The right boundary prefix gives the analogous estimate for
\(Q_{\xi,M_z}^{r,b}\), because \(z[M_z]<\iright^{(M_z)}\) leaves a pile
of size at least \(\chi N\) to its right. Therefore, whenever both
boundary sides have genuine crossings,
\[
    \mathrm{Fluc}_{\rm left}^{r,b}(\xi;y)
    \wedge
    \mathrm{Fluc}_{\rm right}^{r,b}(\xi;z)
    \ge
    N^{c_L^{r,b}(\xi)/2-o(1)}.
\]
If a retained pair is terminal, the stable family is already the singleton
\(\{y\vee z\}\). Such pairs are not used in the anti-concentration cases
below; their contribution is controlled by the terminal/tail estimate in
the numerical and edge-density lemmas.

For a \jscomment{->non-singleton?} first-crossing \(x\), the parent condition gives
\[
    c_L^{r,b}(\mathcal Fx)-c_F^{r,b}(v;y,z)\ge\delta .
\]
Applying the first estimate to \(\mathcal Fx\) and using
\(Q_{\mathcal Fx}\ge0\) gives the desired upper bound for
\(S_{\mathcal Fx}\). Passing from \(\mathcal Fx\) to \(x\) changes \(S\)
by \(O((\log N)^{-1})\) if the final digit is good and by \(0\) if it is
exceptional. This proves the displayed parent-crossing estimate.

Finally fix \(\theta\in[3,4]\). For each \((q,i)\), relative entropy
applied to the distribution
\(c_\ell^{(q,i)}/c_{\rm tot}^{(q,i)}\) and to weights proportional to
\((p_\ell^{(i)})^\theta\) gives, with the convention that zero-count
terms are omitted,
\[
    \sum_\ell c_\ell^{(q,i)}
        \log\frac{c_{\rm tot}^{(q,i)}}{c_\ell^{(q,i)}}
    \le
    \theta\sum_\ell c_\ell^{(q,i)}\log\frac1{p_\ell^{(i)}}
    +c_{\rm tot}^{(q,i)}
        \log\sum_\ell (p_\ell^{(i)})^\theta .
\]
\jscomment{Write the general relative entropy inequality as a lemma before?}
Equivalently,
\[
    H_x\le
    \theta S_x-\sum_{q,i}c_{\rm tot}^{(q,i)}\psi_{\bp^{(i)}}(\theta).
\]
The total number of \((q,i)\)-positions between \(s\) and \(u\) differs
from \(c_{\rm tot}^{(q,i)}\log N\) only by the exceptional positions.
After profiles with \(R_x\le\alpha_\chi\) are retained, interval
blocking and the \(\chi\)-regular representative replacement give
\[
    \sum_{q,i}c_{\rm tot}^{(q,i)}\psi_{\bp^{(i)}}(\theta)
    =
    B_\theta(s,u)+O(\Delta+\chi+\rho+\varphi+\alpha_\chi)+o(1),
\]
uniformly for \(\theta\in[3,4]\). This yields the entropy estimate.
\end{proof}

\begin{lem}[Bad-profile reduction]
For every fixed target \(\mathsf C_0\), the admissibility truncation above
and the bad-profile truncation can be chosen so that the total
contribution of non-admissible boundary pairs and profiles with
\(R_x>\alpha_\chi\) is \(N^{-\mathsf C_0}\). Empty boundary pairs
contribute zero, and terminal retained pairs are kept as singleton
classes. Moreover, for every \(\varepsilon>0\), after choosing \(\chi\)
small and then \(N\) large, the number of remaining
boundary-pair/profile classes is at most \(N^\varepsilon\).
\end{lem}

\jscomment{1.This lemma should be stated in a more rigorous way. Try to avoid use vague words and phrases like "the admissibility truncation above and the bad-profile truncation", the word "contribution"/"contributes zero". 2. What is "remaining boundary-pair/profile classes"? Either recall or define it.}

\begin{proof}
Fix a boundary threshold \(\mathsf C_{\rm bdry}\gg \mathsf C_0\), take
the constant \(\mathsf C\) in the admissibility definition to be
\(\mathsf C_{\rm bdry}\), then choose
\(\chi\) so small that all \jscomment{-->entropy factor below?} are dominated by powers
of \(\chi\), and only then take \(N\to\infty\). Put
\[
    \alpha_\chi=(\log(1/\chi))^{-1/2}.
\]
For any prefix family \(\mathcal V\), write \(G_{\mathcal V}\) for the
subgraph generated by rows in \(B_{\mathcal V}\). Since
\[
    \sum_m \Pbb((m,m+1)\in E(G_{\mathcal V}))^2
    \le
    \sum_m \Pbb((m,m+1)\in E(G_{\mathcal V}))
    =
    \E |E(G_{\mathcal V})|
    \le
    \E |\cI(\mathcal V)| ,
\]
it suffices for discarded families to bound their expected interval mass.

\emph{Admissible boundary truncation.}
If a left boundary prefix is excluded, then it contains
\[
    q>\frac{\mathsf C_{\rm bdry}}{\log(1/\chi)}\log N
\]
digits in piles of size \(<\chi N\) before reaching the boundary stop.
For a fixed set of such locations its expected mass has an extra factor at
most \(\chi^q\le N^{-\mathsf C_{\rm bdry}}\). Summing over the possible
locations and digits gives the explicit tail bound
\[
    N\sum_{q>\mathsf C_{\rm bdry}\log N/\log(1/\chi)}
    \binom Kq k^q\chi^q
    \le
    N^{1-\mathsf C_{\rm bdry}+o_\chi(1)} .
\]
\jscomment{Maybe we should explain more on how to obtain this inequality.}
Here \(K=O(\log N)\), and the \(o_\chi(1)\) term tends to \(0\) after
\(\chi\downarrow0\). Choosing \(\mathsf C_{\rm bdry}\) larger than
\(\mathsf C_0+2\) makes the discarded left-boundary contribution at most
\(N^{-\mathsf C_0}\). The right boundary is identical. Empty pairs
\(A_{y,z}^{r,b}=\varnothing\) contribute zero. Terminal retained pairs are
not included in this discard; they remain as singleton profile classes.

\emph{Exceptional profile truncation.}
If \(R_x>\alpha_\chi\), then \(x\) uses more than
\(\alpha_\chi\log N\) exceptional digits, each carrying a pile-size factor
at most \(3\chi+o(1)\) by \(\chi\)-regularity. Hence their total mass is
bounded by
\[
    \sum_{q\ge\alpha_\chi\log N}\binom Kq k^q(4\chi)^q .
\]
The number of choices of exceptional locations and digits alone \jscomment{replace this with math expression} is
\[
    \exp\{O(\alpha_\chi\log(1/\alpha_\chi)\log N)\}.
\]
\jscomment{Also explain more on how to obtain this. Is the fact $\alpha_{\chi} \to 0$ used?}
Since
\(\alpha_\chi\log(1/\chi)=\sqrt{\log(1/\chi)}\to\infty\), the small-pile
loss dominates the entropy factor, so the mass and hence the edge
contribution of these profiles is \(N^{-\mathsf C_0}\) after decreasing
\(\chi\) \jscomment{Maybe write as "at most" is better, since the first truncation also requires $\chi$ small enough.}.

\jscomment{"Empty boundary pairs contribute zero, and terminal retained pairs are kept as singleton classes." Is the claim proved here or previously?}

After both truncations, the boundary data and the exceptional set have
\jscomment{at most?} \(N^\varepsilon\) possibilities for any prescribed \(\varepsilon>0\),
after \(\chi\) is chosen small\jscomment{Why?}. Each coarse count
\(c_\ell^{(q,i)}\) has only \(O(\log N)\) possible values. Hence the
remaining boundary-pair/profile classes are at most \(N^\varepsilon\)
for all large \(N\).

%\jscomment{I am totally unsure whether the arguments in the paragraph is correct, because I am not fully sure about the meaning of notation "boundary data", "exceptional set" and "remaining boundary-pair/profile classes" here. Try to first give a rigorous definition and then a rigorous proof. }
\end{proof}

\begin{lem}[Hypergeometric estimates]\label{lem:hypergeo-input}
Let \(Y\sim\operatorname{Hyp}(n_1,n,m)\), the size of the intersection
of a uniformly chosen \(n_1\)-set and a uniformly chosen \(m\)-set in an
\(n\)-point population. Put
\[
    \mu=\E Y=\frac{n_1m}{n},
    \qquad
    \sigma^2=
    \frac{n_1m(n-n_1)(n-m)}{n^2(n-1)} .
\]
There are absolute constants \(c,C>0\) such that, for every \(t\ge0\),
\[
    \Pbb(|Y-\mu|\ge t)
    \le
    2\exp\left\{-c\min\left(\frac{t^2}{\mu+t},t\right)\right\},
\]
and, whenever \(\sigma>0\),
\[
    \sup_q \Pbb(Y=q)\le \frac{C}{\sigma}.
\]
\end{lem}

\jscomment{Did we prove this in our universality paper?}

\begin{proof}
The concentration bound is \jscomment{-->Hoeffding's inequality for sampling without
replacement?}. For the point-mass bound, write the hypergeometric mass
function explicitly as a ratio of binomial coefficients. Stirling's
formula gives \(O(\sigma^{-1})\) uniformly in the central window
\(|q-\mu|=O(\sigma)\). The mass function is log-concave in \(q\), so the
maximum is attained in this window up to an absolute constant factor.
\end{proof}

\begin{lem}[Ambient conditioning transfer]
Fix an admissible nonterminal retained boundary pair \((y,z)\), and let
\(\xi\in A_{y,z}^{r,b}\). On an event of probability
\(1-\exp(-N^{\delta-o(1)})\), uniformly over all prefix families appearing
in the endpoint decompositions below, the random counts inside the ambient
block differ from their conditional means by at most a factor
\(N^{O(\delta)+o(1)}\) \jscomment{Is the constant before $\delta$ important in the following proof?} whenever the mean is at least
\(N^{\delta-o(1)}\), and are at most \(N^{O(\delta)+o(1)}\) otherwise.
\jscomment{1. Decompose two cases into two separate sentences. the current version is a bit hard to read. 2. Recall or explain the definition of the random counts.}
Consequently, if \(a,b\le M_{y\vee z}\) realize the left and right
fluctuation scales,
\[
    \E|\cI(P_{\xi,a}^{r,b})|,
    \E|\cI(Q_{\xi,b}^{r,b})|
    \ge N^{2c_F^{r,b}(\xi;y,z)-o(1)},
\]
then every conditional hypergeometric variable \(Y\) constructed from
these two pieces in the anti-concentration proof satisfies
\[
    \E(Y\mid\mathcal G)\ge
    N^{2c_F^{r,b}(\xi;y,z)-O(\delta)-o(1)}
\]
and both complementary population fractions in its variance formula are
at least \(N^{-O(\delta)-o(1)}\).
\jscomment{What is the precise definition of "hypergeometric variable \(Y\) constructed from
these two pieces in the anti-concentration proof(which is below this lemma)", "complementary population fractions" and "its variance formula"?}
\end{lem}

\begin{proof}
All sets used in the endpoint decompositions \jscomment{Recall or explain what do you mean by "endpoint decomposition"?} are prefix cylinders or
finite intersections of prefix cylinders with one column left unrevealed.
There are only \(N^{o(1)}\) such sets after the admissibility and
profile truncations \jscomment{Why? Does this follow from Lemma 5?}. Revealing columns successively \jscomment{Recall the matrix representation here may be helpful.}, each count evolves by
a hypergeometric thinning with deterministic conditional mean \jscomment{"deterministic" in which sense?}. Applying
Lemma~\ref{lem:hypergeo-input} at every reveal step and union bounding over
the \(O(\log N)\) levels gives the stated transfer of means to realized
ambient counts.

\jscomment{1. Did we prove something that can be applied here in our universality paper? 2. If not, I think this paragraph should be rewritten. The proof is too compressed. And this may be the first technical part.}

The nonterminal boundary condition gives a genuine crossing on both sides:
the left crossing has at least \(\chi N\) deterministic mass to its left,
and the right crossing has at least \(\chi N\) deterministic mass to its
right.
\jscomment{What is the precise meaning of "deterministic mass" here. In one shuffle or after multiple shuffles?}
After conditioning inside \(B_w\), the preceding concentration
transfer preserves these side masses up to \(N^{O(\delta)+o(1)}\) factors\jscomment{"factors" in what sense?}.
Thus the fluctuation-maximizing left and right pieces cannot exhaust the
conditional population. 
\jscomment{What is the precise meaning of this sentence? Explain in a math way.}
The displayed lower bound on \(\E(Y\mid\mathcal G)\)
is the transferred version of the defining lower bounds for
\(\E|\cI(P_{\xi,a}^{r,b})|\) and \(\E|\cI(Q_{\xi,b}^{r,b})|\).
\end{proof}

\jscomment{I suggest the proof above should be rewritten in a rigorous way. Remember to avoid using any new concept without explaining its meaning rigorously in math sense before.}

\begin{lem}[Ambient concentration and anti-concentration]
Let \((y,z)\) be a nonterminal retained fluctuation-regular pair and let
\(x\in A_{y,z}^{r,b}\). Put \(M_\vee=|y\vee z|\). Write
\[
    \cI(B_x)=[\iota_x,\tau_x)
\]
for the random lexicographic interval occupied by \(B_x\), and let
\(\mathcal U_x\) be the internal information consisting of the rows whose
strings have prefix \(x\), together with their continuation strings after
time \(M_x\). On an event \(\mathcal S_x\) with
\[
    \Pbb(\mathcal S_x^c)\le \exp(-N^{\delta-o(1)}),
\]
the following estimates hold uniformly over the remaining retained
profiles:
\[
    |\cI(B_x)|\le N^{c_L^{r,b}(x)+O(\delta)+o(1)}
    \qquad\text{if }c_L^{r,b}(x)\ge 2\delta,
\]
\[
    |\cI(B_x)|\le N^{O(\delta)+o(1)}
    \qquad\text{if }c_L^{r,b}(x)<2\delta,
\]
and
\[
    \max_j \Pbb(\iota_x=j,\mathcal S_x\mid\mathcal U_x)
    +
    \max_j \Pbb(\tau_x=j,\mathcal S_x\mid\mathcal U_x)
    \le
    N^{-M_x^*+O(\delta)+o(1)},
\]
where
\[
    M_x^*=\max(c_F^{r,b}(x;y,z),\delta).
\]
Moreover, there are deterministic centers
\(\bar\iota_x=\E\iota_x\) and \(\bar\tau_x=\E\tau_x\) such that, on
\(\mathcal S_x\),
\[
    |\iota_x-\bar\iota_x|+|\tau_x-\bar\tau_x|
    \le
    N^{M_x^*+O(\delta)+o(1)} .
\]
\end{lem}

\jscomment{Some definition can be moved outside the statement of the lemma. A more rigorous definition on $\mathcal U_x$ may be necessary.}

\begin{proof}
The size estimates are the ambient-block forms of the prefix
concentration estimates. Revealing the rows inside \(B_w\) one shuffle
time at a time, the number of rows carrying any fixed legal prefix evolves
by successive hypergeometric thinnings. The concentration part of
Lemma~\ref{lem:hypergeo-input}, iterated over \(O(\log N)\) levels, gives
concentration around the conditional mean; the same bound with
\(t=N^{O(\delta)}\) gives the small-mean estimate. The conditioning on
the ambient block changes only the initial population from \(N\) to
\(|\cI(B_w)|=N^{b+o(1)}\), and all exponents above are measured relative
to this ambient expected size through \(c_L^{r,b}\).

The endpoint-window estimate follows from the same decomposition. For
the left endpoint,
\[
    \iota_x=
    C_w+\sum_{M_w<i\le M_x}|\cI(P_{x,i}^{r,b})|,
\]
where \(C_w\) is the deterministic left endpoint of the ambient block
after conditioning on \(B_w\). Each summand has mean at most
\(N^{2c_F^{r,b}(x;y,z)+o(1)}\): for \(i\le M_\vee\) this is the
definition of the left fluctuation, while for \(i>M_\vee\) the summand is
contained in \(B_{y\vee z}\), whose expected size is at most
\(N^{2c_F^{r,b}(x;y,z)+o(1)}\) by the fluctuation comparison. \jscomment{What is "fluctuation comparison"? Please give a reference.}
The same
argument applies to the right-endpoint decomposition.
Applying Lemma~\ref{lem:hypergeo-input} and taking a union bound over the
\(O(\log N)\) summands gives
\[
    |\iota_x-\bar\iota_x|
    \le N^{M_x^*+O(\delta)+o(1)}
\]
on an event of probability \(1-\exp(-N^{\delta-o(1)})\). Since
\(\tau_x=\iota_x+|\cI(B_x)|\) and the size estimates give the same
deviation scale for \(|\cI(B_x)|\), the displayed window for \(\tau_x\)
follows after intersecting the corresponding good events.
\jscomment{In the proof, the fluctuation of $C_{w}$ has not been taken into consideration. So what one can obtain solely from the arguments above is the result obtained by replacing $\bar \iota_x$ with $\E[\iota_x|B_w]$.}

It remains to identify a hypergeometric variable with variance at least
\(N^{2M_x^*-O(\delta)-o(1)}\)\jscomment{What are you proving and why this is enough? Referring to the results you have used: the distribution and the anti-concentration result you have used.}. Assume first that
\(c_F^{r,b}(x;y,z)\ge\delta-(\log N)^{-1/2}\). Work inside the ambient
population \(B_w\); all row sets below are implicitly intersected with
the rows extending \(w\). Write \(N_w\) for this ambient population size
and, in this paragraph only, write \(n_\ell^{(j)}\) for the corresponding
ambient count in column \(j\).\jscomment{Use another notation for \(n_\ell^{(j)}\)?} Choose columns \(a,b\le M_\vee\) so that
\[
    \E|\cI(P_{x,a}^{r,b})|
    \ge
    N^{2c_F^{r,b}(x;y,z)-o(1)},\qquad
    \E|\cI(Q_{x,b}^{r,b})|
    \ge
    N^{2c_F^{r,b}(x;y,z)-o(1)},
\]
after interchanging left and right if necessary. If the right fluctuation
is the smaller one, reverse the lexicographic order; this swaps \(P\) and
\(Q\) and leaves the argument unchanged.
\jscomment{The letter $b$ has multiple meaning? (Different meaning in $\mathcal A^{r,b}$)}
\jscomment{Is it legal to interchange left and right, since after this exchanging, $\iota_x$ and $\tau_x$ also exchanges.}

We spell out the left-endpoint construction. In the case \(a>b\), reveal
all columns except column \(b\), and in column \(b\) also reveal the rows
in
\[
    R_1=\{\sigma:\ \sigma[M_w+1:b-1]=x[M_w+1:b-1],\
    \sigma[b+1:a]\ge_{\rm lex} x[b+1:a]\}.
\]
Let \(L_b=\sum_{\ell\le x[b]}n_\ell^{(b)}\), with the count restricted to
the current ambient population, and define
\[
    a_2=|\{\sigma:\sigma[M_w+1:b-1]=x[M_w+1:b-1],\
           \sigma[b+1:a]<_{\rm lex}x[b+1:a]\}|,
\]
\[
    a_3=|\{\sigma\in R_1:\sigma[b]\le x[b]\}|.
\]
The remaining part of the left endpoint is \jscomment{What do you mean by "The remaining part of the left endpoint"?}
\[
    Y_1=
    |\{\sigma:\sigma[M_w+1:b-1]=x[M_w+1:b-1],\
       \sigma[b]\le x[b],\ \sigma[b+1:a]<_{\rm lex}x[b+1:a]\}|,
\]
and, conditional on the revealed sigma-field \(\mathcal F_1\supseteq
\mathcal U_x\),
\[
    Y_1\sim\operatorname{Hyp}(L_b-a_3,\ N_w-|R_1|,\ a_2),
    \qquad
    \iota_x=A_1+Y_1,\quad
    \tau_x=A_1+Y_1+|\cI(B_x)|,
\]
for an \(\mathcal F_1\)-measurable shift \(A_1\). The ambient conditioning
transfer lemma, applied on \(\mathcal S_x\)\jscomment{What is the definition of $\mathcal S_x$? And why the following displayed formulae hold It seems that Lemma 7 is applied here, but the statement of Lemma 7 itself is problematic...}, gives
\[
    \E(Y_1\mid\mathcal F_1)\ge N^{2c_F^{r,b}(x;y,z)-O(\delta)-o(1)},
\]
and
\[
    \frac{a_2}{N_w-|R_1|}\le 1-N^{-O(\delta)-o(1)},\qquad
    \frac{L_b-a_3}{N_w-|R_1|}\le 1-N^{-O(\delta)-o(1)}.
\]
With \(n=N_w-|R_1|\), the variance formula gives
\[
\begin{aligned}
    \operatorname{Var}(Y_1\mid\mathcal F_1)
    &=
    \E(Y_1\mid\mathcal F_1)
    \left(1-\frac{L_b-a_3}{n}\right)
    \left(1-\frac{a_2}{n}\right)
    \frac{n}{n-1}\\
    &\ge N^{2c_F^{r,b}(x;y,z)-O(\delta)-o(1)}.
\end{aligned}
\]

The case \(a<b\) is the symmetric reveal at column \(a\). With
\[
    R_2=\{\sigma:\ \sigma[M_w+1:a-1]=x[M_w+1:a-1],\
    \sigma[a+1:b]\le_{\rm lex} x[a+1:b]\},
\]
and \(L_a^-=\sum_{\ell<x[a]}n_\ell^{(a)}\), define \(b_2\) as the number
of rows with the common prefix up to \(a-1\) and
\(\sigma[a+1:b]>_{\rm lex}x[a+1:b]\), and
\(b_3=|\{\sigma\in R_2:\sigma[a]<x[a]\}|\).
Then, conditional on the analogous sigma-field \(\mathcal F_2\),
\[
    Y_2\sim\operatorname{Hyp}(L_a^- - b_3,\ N_w-|R_2|,\ b_2),
    \qquad
    \iota_x=A_2+Y_2,\quad
    \tau_x=A_2+Y_2+|\cI(B_x)|,
\]
\jscomment{mention that $A_2$ is measurable wrt ...}
and the same mean and non-exhaustion estimates \jscomment{What is non-exhaustion estimate?} give
\[
    \E(Y_2\mid\mathcal F_2)\ge N^{2c_F^{r,b}(x;y,z)-O(\delta)-o(1)},
\]
\[
    1-\frac{L_a^- -b_3}{N_w-|R_2|}
    \ge N^{-O(\delta)-o(1)},\qquad
    1-\frac{b_2}{N_w-|R_2|}
    \ge N^{-O(\delta)-o(1)}.
\]
Therefore the hypergeometric variance formula gives
\[
    \operatorname{Var}(Y_2\mid\mathcal F_2)
    \ge N^{2c_F^{r,b}(x;y,z)-O(\delta)-o(1)}.
\]

Finally, if \(a=b\), reveal all columns except \(a\), and in column \(a\)
reveal the rows with digit \(x[a]\). Let \(\mathcal F_3\supseteq
\mathcal U_x\) be the resulting sigma-field. If \(d_2\) is the number of
rows with the common prefix up to \(a-1\) and digit different from
\(x[a]\), then
\[
    Y_3\sim
    \operatorname{Hyp}(L_a^-,\ N_w-n_{x[a]}^{(a)},\ d_2),
    \qquad
    \iota_x=A_3+Y_3,\quad
    \tau_x=A_3+Y_3+|\cI(B_x)|.
\]
The ambient conditioning transfer lemma gives the needed side mass on both
sides of the crossing digit, so
\[
    \E(Y_3\mid\mathcal F_3)\ge N^{2c_F^{r,b}(x;y,z)-O(\delta)-o(1)},
\]
\[
    1-\frac{L_a^-}{N_w-n_{x[a]}^{(a)}}\ge N^{-O(\delta)-o(1)},
    \qquad
    1-\frac{d_2}{N_w-n_{x[a]}^{(a)}}\ge N^{-O(\delta)-o(1)}.
\]
The variance formula gives
\[
    \operatorname{Var}(Y_3\mid\mathcal F_3)
    \ge N^{2c_F^{r,b}(x;y,z)-O(\delta)-o(1)}.
\]
Thus all three cases have conditional variance at least this large.

For every realization of the finer revealed sigma-field inside
\(\mathcal S_x\), the relevant hypergeometric variable has variance at
least \(N^{2c_F^{r,b}(x;y,z)-O(\delta)-o(1)}\). Applying the
anti-concentration part of Lemma~\ref{lem:hypergeo-input} for each such
realization and then taking conditional expectation back to
\(\mathcal U_x\) gives
\[
    \sup_j\Pbb(\iota_x=j,\mathcal S_x\mid\mathcal U_x)
    \le
    N^{-c_F^{r,b}(x;y,z)+O(\delta)+o(1)}.
\]
When \(c_F^{r,b}(x;y,z)<\delta\), the trivial bound
\(\max_j\Pbb(\iota_x=j\mid\mathcal U_x)\le1\) is
\(N^{-\max(c_F^{r,b}(x;y,z),\delta)+O(\delta)}\), after increasing the
implicit constant in \(O(\delta)\). The same argument applies to
\(\tau_x\), using the right endpoint decomposition. This proves the
anti-concentration estimate.
\end{proof}

\begin{lem}[Per-edge stable localization]
For each \(m\in\{1,\ldots,N-1\}\) and each retained admissible boundary
pair \((y,z)\), define
\[
    A_{y,z,m}^{\stable,r,b}
    =
    \left\{
    x\in A_{y,z}^{\stable,r,b}:
    \Pbb(m\in\cI(B_x))\ge \exp(-N^\delta)
    \right\}.
\]
Then \(|A_{y,z,m}^{\stable,r,b}|\le 2Kk\). Moreover,
\[
    \Pbb\left[
    m\notin
    \bigcup_{(y,z)}\ \bigcup_{x\in A_{y,z,m}^{\stable,r,b}}
    \cI(B_x)
    \right]
    \le \exp(-N^{\delta-o(1)}),
\]
where the union is over retained admissible boundary pairs. If
\((m,m+1)\in E(G)\), then the two endpoint rows have the same full string,
so any stable block containing \(m\) also contains the edge.
\end{lem}

\begin{proof}
The global boundary decomposition and the stable partition show that
retained admissible stable intervals partition the retained part of the
ambient block. 
\jscomment{Recall or explain what is "global boundary decomposition" and "stable partition". Please cite something if you have explain this phrase before.}
Intervals omitted by the boundary and profile truncations
have total edge contribution \(N^{-\mathsf C_0}\) by the bad-profile
reduction and are excluded here.
\jscomment{If you are applying Lemma 5, please write explicitly that you have applied this, and explain rigorously how to apply this lemma.}

For a fixed retained pair, discard from \(A_{y,z}^{\stable,r,b}\) all
blocks whose interval contains \(m\) with probability below
\(\exp(-N^\delta)\). A union bound over the at most \(k^K=N^{O(1)}\)
stable prefixes gives the displayed failure probability. It remains to
bound the number of non-discarded blocks. Map every non-singleton
first-crossing block to its parent. If three different parents for the
same boundary pair could contain \(m\) with probability at least
\(\exp(-N^\delta)\), their deterministic centers would be ordered from
left to right; for the middle one, the endpoint-window estimate from the
ambient concentration lemma \jscomment{Which lemma? Cite it as Lemma x.x.} would force one of the outer events to have
probability at most \(\exp(-N^{\delta-o(1)})\), a contradiction for large
\(N\). 
\jscomment{Explain why, or at least give a reference for similar arguments. You should at least use the definition of $A^{\mathrm{stable,r,b}}_{y,z}$ somewhere in the proof of this lemma.}
Thus there are at most two parents, and each has at most \(k\)
children. Terminal singleton pairs contribute one block. The harmless
factor \(K\) covers the possibility of first crossings at different
lengths before applying the parent map uniformly.
\end{proof}

\begin{lem}[Edge-density bounds]
Let \((y,z)\) be a nonterminal retained fluctuation-regular pair and let
\(x\in A_{y,z}^{\stable,r,b}\setminus\Lambda_{y,z}^{r,b}\). The
following bounds hold:
\[
    \sum_{m=1}^{N-1}\Pbb((m,m+1)\in E(G_{B_x}))^2
    \le N^{4c_L^{r,b}(x)-c_F^{r,b}(x;y,z)+2c_E(x)+O(\delta)+o(1)}
\]
in the large-block case \(c_L^{r,b}(x)\ge2\delta\);
\[
    \sum_{m=1}^{N-1}\Pbb((m,m+1)\in E(G_{B_x}))^2
    \le N^{-c_F^{r,b}(x;y,z)+2c_E(x)+O(\delta)+o(1)}
\]
in the small non-singleton case
\(c_L^{r,b}(x)<2\delta\),
\(c_F^{r,b}(x;y,z)\ge\delta-(\log N)^{-1/2}\), and
\(A_{y,z}^{\stable,r,b}\ne\{y\vee z\}\). If
\(c_L^{r,b}(x)<2\delta\) and
\(c_F^{r,b}(x;y,z)<\delta-(\log N)^{-1/2}\), then \(x=y\vee z\) and
\(A_{y,z}^{\stable,r,b}=\{y\vee z\}\). If
\(c_L^{r,b}(x)<2\delta\) and
\(A_{y,z}^{\stable,r,b}=\{y\vee z\}\), then
\[
    \sum_{m=1}^{N-1}\Pbb((m,m+1)\in E(G_{B_x}))^2
    \le N^{2c_L^{r,b}(x)+2c_E(x)+O(\delta)+o(1)}.
\]
\end{lem}

\begin{proof}
For any prefix \(x\) of absolute length \(M_x\),
\[
    \E\bigl(|E(G_{B_x})|\,\big|\,|\cI(B_x)|\bigr)
    \le
    |\cI(B_x)|^2
    \prod_{M_x<j\le K}\sum_\ell
    \frac{n_\ell^{(j)}(n_\ell^{(j)}-1)}{N(N-1)}
    \le |\cI(B_x)|^2N^{c_E(x)+o(1)} .
\]
\jscomment{Maybe add a reference for this displayed formula is better.}

In all cases below we work on the event \(\mathcal S_x\) from the
ambient concentration and anti-concentration lemma. Its complement has
probability \(\exp(-N^{\delta-o(1)})\), so its contribution to the
edge-density sums is absorbed in the displayed \(O(\delta)+o(1)\)
exponents.
Conditional on \(\mathcal U_x\), the internal graph \(G_{B_x}\) is fixed
up to translation by \(\iota_x\). For a fixed \(m\), each internal edge
of this fixed graph is placed at \((m,m+1)\) for at most one value of
\(\iota_x\). Therefore
\[
    \Pbb((m,m+1)\in E(G_{B_x}),\mathcal S_x\mid\mathcal U_x)
    \le
    |E(G_{B_x})|\,
    N^{-M_x^*+O(\delta)+o(1)} .
\]
Averaging this inequality over \(\mathcal U_x\) and using the preceding
conditional expectation bound gives the pointwise estimates \jscomment{Which estimate?} below. The
endpoint-window estimate shows that
all non-negligible \(m\)'s lie in a deterministic interval of size
\[
    N^{c_L^{r,b}(x)+O(\delta)+o(1)}
    +
    N^{M_x^*+O(\delta)+o(1)} .
\]
\jscomment{How did you get this? Did you use Lemma 9 here? If so write this explicitly.}

In the large-block case, the fluctuation comparison \jscomment{What is fluctuation comparison?} and
\(c_L^{r,b}(x)\ge2\delta\) imply
\(c_F^{r,b}(x;y,z)\ge\delta-o(1)\), so \(M_x^*=c_F^{r,b}(x;y,z)+o(1)\).
The size concentration and the internal edge estimate give
\[
    \E\bigl(|E(G_{B_x})|\mathbf 1_{\mathcal S_x}\bigr)
    \le N^{2c_L^{r,b}(x)+c_E(x)+O(\delta)+o(1)}.
\]
Stability \jscomment{of what? Why it gives the following inequality? Where this inequality is applied and what is the meaning of LHS(refer to some equation)?} gives
\[
    N^{c_L^{r,b}(x)+o(1)}+N^{c_F^{r,b}(x;y,z)+o(1)}
    \le N^{c_F^{r,b}(x;y,z)+O(\delta)+o(1)} .
\]
Thus the pointwise edge probability is at most \jscomment{What do you mean by "pointwise edge probability"? Use more math notation.}
\[
    N^{2c_L^{r,b}(x)-c_F^{r,b}(x;y,z)+c_E(x)+O(\delta)+o(1)}.
\]
Squaring and summing over the above window proves the large-block bound.
\jscomment{I remember in previous paper there is a Cauchy inequality applied here. Be more careful here. }

In the small non-singleton case, the extra hypothesis
\(c_F^{r,b}(x;y,z)\ge\delta-(\log N)^{-1/2}\) is exactly what is needed
to use the same anti-concentration scale. The ambient small-block
concentration bound and the internal edge estimate give
\[
    \E\bigl(|E(G_{B_x})|\mathbf 1_{\mathcal S_x}\bigr)
    \le N^{c_E(x)+O(\delta)+o(1)},
\]
and hence the pointwise edge probability is
\[
    N^{-c_F^{r,b}(x;y,z)+c_E(x)+O(\delta)+o(1)}.
\]
The endpoint-window estimate and \(c_L^{r,b}(x)<2\delta\) allow us to
sum its square over a window of size
\(N^{c_F^{r,b}(x;y,z)+O(\delta)+o(1)}\). This gives the second bound.

If \(c_L^{r,b}(x)<2\delta\) and
\(c_F^{r,b}(x;y,z)<\delta-(\log N)^{-1/2}\), then the fluctuation
comparison at \(y\vee z\) gives
\[
    c_L^{r,b}(y\vee z)-c_F^{r,b}(y\vee z;y,z)<\delta
\]
for \(N\) large and the small parameters fixed. Hence the stable family
is the singleton \(\{y\vee z\}\), and \(x=y\vee z\).
\jscomment{Explain on the proof of the displayed equation above briefly.}

For a small singleton stable block, use the ambient anti-concentration
lemma at the scale
\[
    M_*=\max(c_F^{r,b}(x;y,z),\delta).
\]
The endpoint-window estimate gives a relevant window of size
\[
    N^{c_L^{r,b}(x)+O(\delta)+o(1)}
    +
    N^{M_*+O(\delta)+o(1)} ,
\]
and the small-mean concentration bound gives the second-moment estimate
\[
    \E|\cI(B_x)|^2\le N^{c_L^{r,b}(x)+O(\delta)+o(1)}.
\]
Indeed, on \(\mathcal S_x\) we have
\(|\cI(B_x)|\le N^{O(\delta)+o(1)}\), so
\(|\cI(B_x)|^2\le N^{O(\delta)+o(1)}|\cI(B_x)|\); the complement of
\(\mathcal S_x\) has exponentially small probability.
Consequently the pointwise probability is bounded by
\[
    N^{c_L^{r,b}(x)-M_*+c_E(x)+O(\delta)+o(1)}.
\]
Multiplying the window size by the square of this pointwise bound gives \jscomment{Again, why you need to estimate LHS in the formula below?}
\[
\begin{aligned}
&\left(
N^{c_L^{r,b}(x)+O(\delta)+o(1)}
    +
    N^{M_*+O(\delta)+o(1)}
\right)\\
&\quad \times
N^{2c_L^{r,b}(x)-2M_*+2c_E(x)+O(\delta)+o(1)}
    \le
    N^{2c_L^{r,b}(x)+2c_E(x)+O(\delta)+o(1)}.
\end{aligned}
\]
Here we used the singleton stability inequality
\(c_L^{r,b}(x)-c_F^{r,b}(x;y,z)<\delta\), which implies
\[
    \max\{c_L^{r,b}(x),M_*\}-2M_*
    \le O(\delta).
\]
This proves the singleton bound.
\end{proof}

\begin{lem}[Terminal singleton bound]
Let \((y,z)\) be a terminal retained admissible pair and
\(x=y\vee z\), so \(A_{y,z}^{\stable,r,b}=\{x\}\) and \(M_x=K\). Then
\[
    \sum_{m=1}^{N-1}\Pbb((m,m+1)\in E(G_{B_x}))^2
    \le N^{c_L^{r,b}(x)+c_E(x)+O(\delta)+o(1)} .
\]
\end{lem}

\begin{proof}
Since \(M_x=K\), the tail factor is \(N^{c_E(x)}=1\). As above,
\[
    \sum_m\Pbb((m,m+1)\in E(G_{B_x}))^2
    \le \E |E(G_{B_x})| .
\]
The small-mean concentration estimate gives
\(\E|\cI(B_x)|^2\le N^{c_L^{r,b}(x)+O(\delta)+o(1)}\), and
\(|E(G_{B_x})|\le |\cI(B_x)|^2\). This proves the claim, with the
displayed \(c_E(x)\) included for notational compatibility.
\end{proof}

\begin{lem}[Numerical lemma]
Assume \(\mathsf{FM}(r,b,T,\eta)\) \jscomment{For all pairs $(r,b)$ or one fixed pair $(r,b)$?}. After choosing
\(\Delta,\chi,\rho,\varphi\) and the stability threshold sufficiently
small, there is \(\delta_1>0\) such that every stable profile remaining
after the admissibility and exceptional-profile truncations satisfies
\[
    4c_L^{r,b}(x)-c_F^{r,b}(x;y,z)+2c_E(x)+H_x\le-\delta_1
\]
in the nonterminal large-block case. In the nonterminal small
non-singleton case,
\[
    -c_F^{r,b}(x;y,z)+2c_E(x)+H_x\le-\delta_1.
\]
Every singleton stable profile, including terminal retained pairs,
satisfies
\[
    c_L^{r,b}(y\vee z)+c_E(y\vee z)\le-\delta_1.
\]
\end{lem}

\jscomment{The lemma should be stated for $y$ and $z$ are continuation of $w$, where $w$ is used at the beginning of Definition 3 (involved $r$ and $b$)? Also, the $\delta_1$ should be uniform over $r$ and $b$ at least?}

\begin{proof}
Let \(\delta\) denote the stability threshold in the definition of
\(A_{y,z}^{\stable,r,b}\), and set
\[
    \varepsilon_{\rm det}=O(\Delta+\chi+\rho+\varphi+\alpha_\chi).
\]
Choose \(\Delta,\chi,\rho,\varphi\), and then \(\delta\), so that the
fixed multiples of \(\varepsilon_{\rm det}+\delta\), including the
edge-density \(O(\delta)\) losses used later, are at most \(\eta/2\), and
set \(\delta_1=\eta/4\). Finally take \(N\) large so all \(o(1)\) errors
are below the remaining margins. For nonterminal retained pairs, throughout,
\[
    c_F^{r,b}(x;y,z)=c_F^{r,b}(y\vee z;y,z)
    \qquad (x\in A_{y,z}^{r,b}).
\]

First consider a terminal retained pair. Then
\(A_{y,z}^{\stable,r,b}=\{y\vee z\}\) and \(s=T\). The maximal-pile and
tail estimates \jscomment{refer to equations like (x.x) explicitly} give
\[
    c_L^{r,b}(y\vee z)+c_E(y\vee z)
    \le \beta_{r,b}(T)+\varepsilon_{\rm det}+o(1)
    \le-\eta+\varepsilon_{\rm det}+o(1),
\]
so the singleton conclusion holds for terminal pairs.
\jscomment{Recall that $s=t_{y \vee z}$.}

Now suppose \((y,z)\) is nonterminal. The terminal fallback
\(\Lambda_{y,z}^{r,b}\) does not contribute among the remaining retained
profiles. If \(|x|=K\), then the
maximal-pile estimate gives
\[
    c_L^{r,b}(x)\le \beta_{r,b}(T)+\varepsilon_{\rm det}+o(1)
    \le-\eta+\varepsilon_{\rm det}+o(1).
\]
Using the fluctuation comparison and the equality of \(c_F\) along the
boundary pair,
\[
    c_L^{r,b}(x)-c_F^{r,b}(y\vee z;y,z)
    =
    c_L^{r,b}(x)-c_F^{r,b}(x;y,z)
    \le \frac12c_L^{r,b}(x)+o(1)<\delta .
\]
Thus \(x\notin\Lambda_{y,z}^{r,b}\), and every non-singleton block below
is a first-crossing block.

\medskip
\noindent\emph{Large-block exponent.}
Fix \(\theta\in[3,4]\). By the profile estimates\jscomment{Which Lemma x.x?},
\[
    H_x+2c_E(x)
    \le \theta S_x-B_\theta(s,u)-2B_2(u,T)
    +O(\varepsilon_{\rm det})+o(1),
\]
and
\[
    c_L^{r,b}(x)
    =
    c_L^{r,b}(y\vee z)-S_x-Q_x+O(\varepsilon_{\rm det})+o(1).
\]
Since \(x\) is stable,
\[
    S_x
    \ge
    c_L^{r,b}(y\vee z)-c_F^{r,b}(y\vee z;y,z)
    -Q_x-\delta-O(\varepsilon_{\rm det})-o(1).
\]
Because \(\theta-4\le0\), this lower bound on \(S_x\) gives an upper
bound when substituted into \((\theta-4)S_x\). Using \(Q_x\ge0\) and
\(c_F^{r,b}(y\vee z;y,z)\ge\frac12c_L^{r,b}(y\vee z)-o(1)\),
\[
\begin{aligned}
    4c_L^{r,b}(x)-c_F^{r,b}(x;y,z)+\theta S_x
    &\le
    \theta c_L^{r,b}(y\vee z)
    +(3-\theta)c_F^{r,b}(y\vee z;y,z)
    +O(\varepsilon_{\rm det}+\delta)+o(1)\\
    &\le
    \frac{\theta+3}{2}\beta_{r,b}(s)
    +O(\varepsilon_{\rm det}+\delta)+o(1).
\end{aligned}
\]
Therefore
\[
    4c_L^{r,b}(x)-c_F^{r,b}(x;y,z)+2c_E(x)+H_x
    \le
    \frac{\theta+3}{2}\beta_{r,b}(s)
    -B_\theta(s,u)-2B_2(u,T)
    +O(\varepsilon_{\rm det}+\delta)+o(1).
\]
Taking \(\theta\) near the infimum in \(\mathcal N_1(r,s,u;b)\) gives
the large-block bound.

\medskip
\noindent\emph{Small non-singleton exponent.}
In the non-singleton first-crossing case, the parent condition gives
\[
    S_x
    \le
    c_L^{r,b}(y\vee z)-c_F^{r,b}(y\vee z;y,z)
    +O(\varepsilon_{\rm det})+o(1).
\]
Since the coefficient of \(c_F\) below is negative, the fluctuation lower
bound gives
\[
\begin{aligned}
    -c_F^{r,b}(x;y,z)+\theta S_x
    &\le
    -c_F^{r,b}(y\vee z;y,z)
    +\theta\big(c_L^{r,b}(y\vee z)-c_F^{r,b}(y\vee z;y,z)\big)
    +O(\varepsilon_{\rm det})+o(1)\\
    &\le
    \frac{\theta-1}{2}\beta_{r,b}(s)
    +O(\varepsilon_{\rm det})+o(1).
\end{aligned}
\]
Combining with entropy and tail estimates and optimizing over
\(\theta\in[3,4]\) gives the small non-singleton bound from
\(\mathcal N_2(r,s,u;b)\).

\medskip
\noindent\emph{Singleton exponent.}
If \(A_{y,z}^{\stable,r,b}=\{y\vee z\}\), then \(u=s\). The profile
estimate for \(c_E\) and the maximal-pile bound give
\[
    c_L^{r,b}(y\vee z)+c_E(y\vee z)
    \le \beta_{r,b}(s)-B_2(s,T)
    +O(\varepsilon_{\rm det})+o(1)
    =
    \mathcal N_3(r,s;b)+O(\varepsilon_{\rm det})+o(1).
\]
The condition \(\mathsf{FM}(r,b,T,\eta)\) makes the right-hand side
negative with margin. This completes the proof.
\end{proof}

\begin{prop}[First moment]
Fix \(\eta>0\), then choose
\(\chi,\Delta,\rho,\varphi,\delta\) according to the parameter hierarchy
above, with the almost-like condition imposed for the chosen
\(\chi\)-regular partition. Fix an ambient block \(B_w\) with parameters
\((r,b)\). Assume the deterministic tail pile sequence
\((\bn^{(j)})_{M_w<j\le K}\), viewed at macroscopic times
\([r,T]\), is almost-\(\mu_\cdot\)-like, satisfies the local
\((\chi,\gamma,\rho)\)-good-shuffle condition on \([r,T]\), and satisfies
\(\mathsf{FM}(r,b,T,\eta)\). Then two independent shuffle graphs
restricted to the ambient block satisfy
\[
    \E |E(G,G')|\le N^{-\delta_*}.
\]
Here \(E(G,G')\) is the set of common adjacent edges,
\(\delta_*=\delta_*(\eta,\gamma,k,T)>0\), and the whole deck is
\(r=0,b=1\).
\end{prop}

\jscomment{The statement is only stated in the $(r,b)=(0,1)$ case, but in the exponential moments part we may need all $(r,b)$ case. If fix $w$ with parameter $(r,b)$, then in principle only results within the block $B_w$ can be obtained.}

\begin{proof}
\[
    \E |E(G,G')|
    =
    \sum_{m=1}^{N-1}\Pbb((m,m+1)\in E(G))^2 .
\]
Fix \(\mathsf C_0\) large. The bad-profile reduction removes
non-admissible boundary pairs and exceptional profiles at cost
\(N^{-\mathsf C_0}\). For the retained part, the per-edge localization
lemma gives, for each \(m\),
\[
    \Pbb((m,m+1)\in E(G))
    \le
    N^{-\mathsf C_0}
    +\sum_{(y,z)}
    \sum_{x\in A_{y,z,m}^{\stable,r,b}}
    \Pbb((m,m+1)\in E(G_{B_x})),
\]
where the sum is over retained admissible pairs. For fixed \(m\), the
number of summands is at most \(N^\varepsilon\) after \(\chi\) is chosen
small, because there are \(N^\varepsilon\) retained boundary/profile
classes and at most \(2Kk=N^{o(1)}\) stable blocks per pair containing
\(m\) with non-negligible probability. Hence AM-GM gives
\[
    \E |E(G,G')|
    \le
    N^{-\mathsf C_0}
    +
    N^\varepsilon
    \max_{\mathcal C}
    \sum_{x\in\mathcal C}
    \sum_{m=1}^{N-1}
    \Pbb((m,m+1)\in E(G_{B_x}))^2 ,
\]
where \(\mathcal C\) ranges over remaining boundary-pair/profile classes.

Fix such a class. For a non-singleton class, the entropy estimate gives
\(|\mathcal C|\le N^{H_x+o(1)}\)\jscomment{explain why}; for singleton classes, \(H_x=0\). The
exceptional locations and digit values have already been included in the
class count through the \(N^\varepsilon\) loss.

First suppose \(c_L^{r,b}(x)\ge2\delta\). This includes possible large
nonterminal singleton classes; they are routed through the large-block
estimate. The edge-density and numerical lemmas give
\[
    N^{H_x}
    N^{4c_L^{r,b}(x)-c_F^{r,b}(x;y,z)+2c_E(x)+O(\delta)+o(1)}
    \le N^{-\delta_1+O(\delta)+o(1)}.
\]
If \(c_L^{r,b}(x)<2\delta\), \(A_{y,z}^{\stable,r,b}\ne\{y\vee z\}\),
and \(c_F^{r,b}(x;y,z)\ge\delta-(\log N)^{-1/2}\), then the small
non-singleton estimates give
\[
    N^{H_x}
    N^{-c_F^{r,b}(x;y,z)+2c_E(x)+O(\delta)+o(1)}
    \le N^{-\delta_1+O(\delta)+o(1)}.
\]
All remaining nonterminal small cases are singleton classes: the
edge-density lemma identifies the low-\(c_F\) small case as singleton,
while high-\(c_F\) small non-singletons were handled above. For such a
small singleton class there is only \(x=y\vee z\), so no entropy
multiplier is needed, and
\[
    N^{2c_L^{r,b}(x)+2c_E(x)+O(\delta)+o(1)}
    \le N^{-2\delta_1+O(\delta)+o(1)}.
\]
Finally, if the class is a terminal singleton, the terminal singleton
bound and the numerical singleton estimate give
\[
    N^{c_L^{r,b}(x)+c_E(x)+O(\delta)+o(1)}
    \le N^{-\delta_1+O(\delta)+o(1)}.
\]
The parameters were chosen so the \(O(\delta)\) and \(o(1)\) losses are
absorbed by the margins. The class-count exponent \(\varepsilon\) is
chosen smaller than the remaining margin after \(\chi\) is fixed, and
\(\mathsf C_0\) is chosen large. This proves the bound for some
\(\delta_*(\eta,\gamma,k,T)>0\).
\end{proof}

\begin{rmk}
The condition above is only an upper-bound first-moment condition. It is not yet a claimed cutoff formula.
\end{rmk}

\end{document}
